\documentclass{article}
\usepackage[utf8]{inputenc}
\usepackage{graphicx}
\usepackage{titling}
\usepackage{float}


\setlength{\parindent}{0pt}
% 1. Set professional margins (1 inch is standard)
\usepackage[margin=1in]{geometry}

% 2. Adjust the title height
\setlength{\droptitle}{-0.5in} 

\title{Labwork 1 - Gradient Descent}
\author{Do Minh Quang - 2540095}
% \date{May 2026}

\begin{document}
\maketitle

\section{Gradient Descent Implementation}
This labwork is for implementing gradient descent algorithm from scratch.

In this labwork, I will work on $f(x) = x^2$ and the goal is to find the local minimum of the function. For the setup, I will experiment the algorithm with the initial value is $x_{0} = 5$. The core update rule of the algorithm is defined as below:
\[x_{n+1} = x_{n} - lr \times f'(x_n)\]
where:
\begin{itemize}
    \item $x_{n}$: is the current value.
    \item $x_{n+1}$: is the update value.
    \item $lr$: is the learning rate.
    \item $f'(x)$: is the first order derivative at $x_n$.
\end{itemize}

As the choice of learning rate has large effect on the performance of the algorithm, I have tried 4 different learning rates, which are 0.001, 0.005, 0.1, 0.5. Below is the convergence speed of each learning rate:

\begin{table}[H]
\centering
\label{tab:converge_iter}
\begin{tabular}{|c|c|}
\hline
\textbf{Learning rate} & \textbf{Converged iteration} \\ \hline
0.001                  & 186901                       \\ \hline
0.005                  & 37231                        \\ \hline
0.1                    & 1677                         \\ \hline
0.5                    & 1                            \\ \hline
\end{tabular}
\end{table}

The table \ref{tab:converge_iter} shows that the convergence speed is highly dependent on the learning rate, with smaller values like 0.001 and 0.005 requiring an large number of iterations. In contrast, $lr = 0.5$ is the most efficient in this case, it achieve the local minimum in a single step. 

\end{document}